% =====================================================================
% NEURAX: An Analytical Compiler for Neural Network Architectures
% Prepared for submission to arXiv (cs.LG)
% Compile: pdflatex main.tex (twice for references)
% =====================================================================
\documentclass[11pt]{article}

% ---- Packages (arXiv-compatible: no exotic dependencies) ----
\usepackage[margin=1in]{geometry}
\usepackage{amsmath,amssymb}
\usepackage{booktabs}
\usepackage{array}
\usepackage{xcolor}
\usepackage{listings}
\usepackage[hidelinks]{hyperref}
\usepackage{titlesec}
\titleformat{\section}{\large\bfseries}{\thesection.}{0.5em}{}
\titleformat{\subsection}{\normalsize\bfseries}{\thesubsection.}{0.5em}{}

% ---- Listings style for Rust code ----
\definecolor{rustcomment}{RGB}{87,96,111}
\definecolor{rustkeyword}{RGB}{0,119,170}
\definecolor{ruststring}{RGB}{163,21,21}
\lstdefinelanguage{RustModern}{
  keywords={let,mut,fn,if,else,while,for,struct,enum,impl,trait,pub,use,mod,as,const,static,match,return,self,true,false,where},
  keywordstyle=\color{rustkeyword}\bfseries,
  comment=[l]{//},
  commentstyle=\color{rustcomment}\itshape,
  string=[b]",
  stringstyle=\color{ruststring},
  sensitive=true,
}
\lstset{
  language=RustModern,
  basicstyle=\ttfamily\footnotesize,
  numbers=left,
  numberstyle=\tiny\color{gray},
  frame=single,
  rulecolor=\color{gray!40},
  breaklines=true,
  xleftmargin=0.25in,
  aboveskip=1em,
  belowskip=1em,
}

% ---- Hyperlinks ----
\urlstyle{same}
\emergencystretch=1em

\title{\textbf{NEURAX: An Analytical Compiler for Neural Network Architectures}}
\author{Martial Fossouo\\
Independent Researcher\\
\texttt{martialwato50@gmail.com}}
\date{August 7, 2026}

\begin{document}
\maketitle

\begin{abstract}
Training modern deep learning models is expensive and unpredictable. Practitioners routinely commit GPU resources to architectures whose memory footprint, training cost, and runtime behavior are unknown until training begins, leading to wasted compute, out-of-memory failures, and slow iteration cycles. We present \textbf{NEURAX}, an analytical compiler for neural network architectures that predicts training cost, peak memory usage, latency, and inference behavior \emph{before} training, in under 50 milliseconds and with zero GPU requirement. NEURAX operates at design time through an eleven-phase intermediate representation (IR) pipeline composed of ten analytical dialects. The system supports ten architecture families---Transformer, CNN, MoE, Diffusion, GNN, RNN, SSM, GAN, Hybrid, and Multimodal---with 680+ configurable blocks and 88 reference templates. Analytical predictions are validated against real training runs with a reported average accuracy of 99\%+ across validated families. NEURAX is implemented in Rust with an MLIR/LLVM 18 backend and is released under the MIT license.
\end{abstract}

\section{Introduction}

\subsection{Motivation}

The cost of training large neural networks has grown superlinearly with model scale. A GPT-3-class model (175B parameters) requires thousands of GPU-days and millions of dollars in compute. Yet the dominant workflow remains trial-and-error: engineers design an architecture, launch training, and discover only afterward whether the model fits in memory, how long training takes, and whether inference is stable.

This workflow has three concrete failure modes:
\begin{enumerate}
  \item \textbf{Memory failures (OOM):} architectures exceed GPU VRAM after days of expensive setup.
  \item \textbf{Cost overruns:} training duration and cloud cost are unknown at launch time.
  \item \textbf{Behavioral surprises:} models exhibit instability or hallucination that only appears at inference time.
\end{enumerate}

\subsection{Contribution}

We introduce \emph{analytical compilation} as a distinct paradigm between training frameworks (PyTorch, TensorFlow), which execute models, and runtime compilers (IREE, XLA), which lower them for execution. NEURAX treats architecture design as a compilation problem: it parses a declarative model description, lowers it through an intermediate representation, and emits a \emph{prediction report}---analogous to a compiler's optimization report---before any compute is committed.

The key insight is that for a fixed architecture and hardware configuration, resource requirements are \emph{analytically computable} from structural properties (layer types, dimensions, sequence length, parallelism strategy). No execution or profiling is required.

\subsection{Positioning}

\begin{table}[ht]
\centering
\caption{Positioning of analytical compilation relative to existing tools.}
\begin{tabular}{@{}l l l l@{}}
\toprule
\textbf{Category} & \textbf{Representative Tools} & \textbf{Operates At} & \textbf{Output} \\
\midrule
Training frameworks & PyTorch, TensorFlow, JAX & Execution time & Trained models \\
Runtime compilers & IREE, XLA, TensorRT & Compile time & Optimized executables \\
\textbf{Analytical compilers} & \textbf{NEURAX} & \textbf{Design time} & \textbf{Prediction reports} \\
\bottomrule
\end{tabular}
\end{table}

\section{System Architecture}

NEURAX is organized as a Rust workspace of nine crates (Figure~\ref{fig:arch}).

\begin{figure}[ht]
\centering
\begin{verbatim}
-------------------------------------------------------------
|                    neurax-cli (CLI)                       |
|                     neurax-service (API)                  |
-------------------------------------------------------------
|                     neurax-core                           |
|            (pipeline orchestration, 11 phases)            |
----------------+--------------+--------------+-------------
|  neurax-ir    | neurax-parser| neurax-      | neurax-     |
| (10-dialect   | (JSON -> AST)| formulas     | hardware-db |
|  IR)          |              | (analytical) | (GPU/CPU DB)|
----------------+--------------+--------------+-------------
|              neurax-mlir (MLIR/LLVM 18 backend)           |
-------------------------------------------------------------
\end{verbatim}
\caption{NEURAX workspace: nine crates spanning parsing, analytical IR, formulas, hardware database, and the MLIR backend.}
\label{fig:arch}
\end{figure}

\subsection{Pipeline Phases}

The analytical pipeline executes eleven phases, each a compiler pass over the IR:
\begin{enumerate}
  \item \textbf{Parse} --- deserialize JSON into a typed \texttt{ModelConfig} AST
  \item \textbf{Architecture} --- validate topology and family-specific parameters
  \item \textbf{Graph} --- construct the computational graph
  \item \textbf{Tensor} --- propagate tensor shapes and dtypes
  \item \textbf{Operator} --- resolve operator-level cost primitives
  \item \textbf{Compute} --- aggregate FLOPs by layer and family
  \item \textbf{Memory} --- compute activation and parameter memory
  \item \textbf{Parallelism} --- evaluate data/tensor/pipeline parallel strategies
  \item \textbf{Hardware} --- apply GPU/CPU specifications
  \item \textbf{Cost} --- convert compute to time, energy, and cloud cost
  \item \textbf{Report + Dynamic Analysis} --- emit report and behavioral metrics (M36--M55)
\end{enumerate}

Phases 8--11 run with parallelism-aware scheduling; dynamic passes (virtual memory, stability, behavioral synthesis) execute concurrently via \texttt{rayon::join}.

\subsection{Analytical IR Dialects}

The IR is organized into ten dialect modules (analogous to MLIR dialects), summarized in Table~\ref{tab:dialects}.

\begin{table}[ht]
\centering
\caption{The ten analytical IR dialects.}
\begin{tabular}{@{}l l@{}}
\toprule
\textbf{Dialect} & \textbf{Responsibility} \\
\midrule
\texttt{Architecture} & Model topology, global parameters \\
\texttt{Tensor} & Shape and dtype propagation \\
\texttt{Operator} & Operator-level cost primitives \\
\texttt{Compute} & FLOPs and arithmetic intensity \\
\texttt{Memory} & Activation + parameter memory \\
\texttt{Parallelism} & Strategy evaluation \\
\texttt{Hardware} & Device specifications \\
\texttt{Cost} & Time, energy, pricing \\
\texttt{Report} & Report generation \\
\texttt{Precision} & Confidence estimation \\
\bottomrule
\end{tabular}
\label{tab:dialects}
\end{table}

The \texttt{neurax-mlir} backend lowers these dialects to MLIR textual form using the \texttt{melior} bindings, targeting LLVM 18.

\section{Analytical Methodology}

\subsection{Cost Model}

For a neural network with sequence length $S$, batch size $B$, hidden dimension $H$, and number of attention heads $h$ (head dimension $d_h = H/h$), the dominant FLOP contributions are computed as closed-form expressions, implemented verbatim in the \texttt{neurax-formulas} crate:

\textbf{Attention} (non-causal). QKV projections contribute $6BSH^2$; attention scores contribute $2BS^2H$; the attention-to-value matmul contributes $2BS^2H$; softmax contributes $5BS^2h$; the output projection contributes $2BSH^2$:
\begin{equation}
\mathrm{FLOPs}_{\mathrm{attn}} = 8BSH^2 + 4BS^2H + 5BS^2h
\label{eq:attn}
\end{equation}

For \emph{causal} attention the quadratic terms are halved ($\times 0.5$). Grouped-query attention (GQA) reduces the K/V projections to $d_h$ dimensions per KV head ($d_{kv}$), giving $2 \cdot 2BSH d_{kv}$ instead of $6BSH^2$ for the projections.

\textbf{MLP.} A two-layer MLP with intermediate dimension $I$ contributes one matmul $H \to I$ and one $I \to H$:
\begin{equation}
\mathrm{FLOPs}_{\mathrm{mlp}} = 2BSHI + 2BSIH = 4BSHI
\label{eq:mlp}
\end{equation}

\textbf{Embedding.} A lookup over vocabulary $V$ and dimension $H$ is treated as a memory-bound operation:
\begin{equation}
\mathrm{FLOPs}_{\mathrm{emb}} = BSH
\label{eq:emb}
\end{equation}

Total forward FLOPs for a transformer with $N$ layers:
\begin{equation}
\mathrm{FLOPs}_{\mathrm{fwd}} = N \cdot (\mathrm{FLOPs}_{\mathrm{attn}} + \mathrm{FLOPs}_{\mathrm{mlp}}) + \mathrm{FLOPs}_{\mathrm{emb}}
\label{eq:fwd}
\end{equation}

The backward pass approximately triples the forward cost, i.e.\ $\mathrm{FLOPs}_{\mathrm{train}} \approx 3 \cdot \mathrm{FLOPs}_{\mathrm{fwd}}$ \cite{kaplan2020scaling}. NEURAX uses family-specific closed-form formulas for each of the ten supported families, implemented in \texttt{neurax-formulas}; representative primitives include \texttt{attention\_flops}, \texttt{gqa\_flops}, \texttt{mlp\_flops}, \texttt{resnet\_bottleneck\_block\_flops}, and \texttt{mbconv\_flops}.

\subsection{Memory Model}

Peak VRAM usage decomposes into:
\begin{equation}
M_{\mathrm{peak}} = M_{\mathrm{params}} + M_{\mathrm{optimizer}} + M_{\mathrm{activations}} + M_{\mathrm{gradients}}
\label{eq:mem}
\end{equation}
where $M_{\mathrm{params}} = P \cdot \mathrm{bytes}$ and optimizer state for AdamW adds $12$ bytes/parameter (2 copies of moments) plus gradients ($4$ bytes/parameter). Activation memory is computed by propagating tensor shapes through the graph with the configured sequence length and batch size.

\subsection{Behavioral Synthesis}

Beyond resource prediction, NEURAX performs \emph{dynamic analysis} to predict inference behavior without execution:
\begin{itemize}
  \item \textbf{Stability analysis} (M36--M39): entropy evolution and attention focus metrics
  \item \textbf{Virtual memory} (M40--M45): virtual memory pressure under concurrency
  \item \textbf{Behavioral synthesis} (M46--M55): hallucination-risk proxy, inference stability, and 22 configurable inference parameters (sampling, context, stress testing)
\end{itemize}
These are analytical proxies derived from architectural properties (e.g., depth, width, normalization placement), not runtime measurements.

\section{Supported Architectures}

NEURAX supports ten model families and 680+ configurable blocks, summarized in Table~\ref{tab:families}.

\begin{table}[ht]
\centering
\caption{Supported architecture families.}
\begin{tabular}{@{}l l l@{}}
\toprule
\textbf{Family} & \textbf{Coverage} & \textbf{Example Templates} \\
\midrule
Transformer & Attention, MLP, embeddings, GQA & GPT-4, LLaMA 3, Mixtral \\
CNN & ResNet, MobileNet, ConvNeXt, DenseNet & ResNet-50, MobileNet \\
MoE & Expert routing, top-k & Mixtral \\
Diffusion & UNet, denoising & Stable Diffusion, SDXL \\
GNN & Graph convolution & GraphSAGE \\
RNN & LSTM, GRU & Seq2Seq \\
SSM & State space, Mamba blocks & Mamba, S4, H3 \\
GAN & Generator/discriminator & DCGAN \\
Hybrid & Multi-architecture (ViT, DiT) & ViT \\
Multimodal & Vision-language & CLIP, LLaVA \\
\bottomrule
\end{tabular}
\label{tab:families}
\end{table}

88 reference templates cover production architectures from GPT-4 to Stable Diffusion.

\section{Validation Methodology}

\subsection{Benchmark Protocol}

Predictions are validated against real training runs using the following protocol:
\begin{enumerate}
  \item \textbf{Select} a model configuration (JSON)
  \item \textbf{Analyze} with NEURAX --- collect predicted metrics
  \item \textbf{Train} the model with standard frameworks
  \item \textbf{Compare} predicted vs.\ measured metrics
  \item \textbf{Report} accuracy per family
\end{enumerate}

\subsection{Validated Metrics}
\begin{itemize}
  \item FLOPs (floating-point operations)
  \item Peak VRAM (GB)
  \item Training time (wall-clock)
  \item Cloud compute cost (USD)
  \item Energy (kWh)
\end{itemize}

\subsection{Results Summary}

\begin{table}[ht]
\centering
\caption{Validation accuracy by model family.}
\begin{tabular}{@{}l c c@{}}
\toprule
\textbf{Model Family} & \textbf{Models Tested} & \textbf{Avg Accuracy} \\
\midrule
Transformer & 10 & 99.2\% \\
CNN & 8 & 99.5\% \\
MoE & 5 & 98.7\% \\
Diffusion & 4 & 98.9\% \\
SSM & 3 & 97.8\% \\
\bottomrule
\end{tabular}
\label{tab:results}
\end{table}

The cross-family average is above 99\% for validated families. Full benchmark data is available at \url{https://rustnew.github.io/NEURAX/benchmarks}, and the validation harness is open source (\texttt{neurax-core/tests/flops\_cross\_validation.rs}, \texttt{real\_world\_comparison.rs}).

\section{Performance}

NEURAX completes a full eleven-phase analysis of 8B-parameter models in \textbf{under 50\,ms} on commodity hardware, with zero GPU requirement. This is enabled by:
\begin{itemize}
  \item Closed-form analytical formulas (no simulation)
  \item Rayon-based parallel pass scheduling
  \item \texttt{ahash}/\texttt{dashmap}-based fast indexing
  \item Small constant-factor overhead per pass
\end{itemize}
The pipeline is fully deterministic: identical input always produces identical output.

\section{Related Work}
\begin{itemize}
  \item \textbf{Training frameworks} (PyTorch, TensorFlow, JAX) execute models but cannot predict resource needs without running them.
  \item \textbf{Runtime compilers} (IREE, XLA, TensorRT, TVM) optimize execution but require a trained model.
  \item \textbf{Memory estimators} (e.g., PyTorch \texttt{torch.cuda} estimates) are framework-bound and operate on live tensors.
  \item \textbf{NAS tools} (Microsoft Archai) search over architectures but do not provide analytical cost prediction at design time.
  \item \textbf{MLIR/LLVM infrastructure} (LLVM 18, \texttt{melior}) provides the compilation substrate NEURAX leverages for its backend.
\end{itemize}
NEURAX is complementary to all of the above: it provides the design-time prediction layer that the ecosystem currently lacks.

\section{Conclusion and Future Work}

We presented NEURAX, an analytical compiler that predicts resource requirements and inference behavior of neural architectures before training, in under 50\,ms, with zero GPU. Its design-time paradigm is complementary to training frameworks and runtime compilers, and its validation results suggest that analytical cost prediction is both feasible and accurate.

Future directions include:
\begin{enumerate}
  \item \textbf{Expanded validation} --- additional families and larger model scales
  \item \textbf{Fine-tuning cost projections} --- LoRA/QLoRA estimates
  \item \textbf{Hugging Face Hub integration} --- predict costs before downloading models
  \item \textbf{Kubernetes/cloud deployment integration}
  \item \textbf{Collaborative multi-user design (CRDT)}
  \item \textbf{Community benchmarks} --- open dataset for third-party validation
\end{enumerate}
NEURAX is open source (MIT) at \url{https://github.com/rustnew/NEURAX}.

\section*{Acknowledgments}
Built with MLIR/LLVM, Rust, and React ecosystems. The analytical formulas draw on established scaling-law literature~\cite{kaplan2020scaling,hoffmann2022training}.

\bibliographystyle{plain}
\begin{thebibliography}{9}

\bibitem{kaplan2020scaling}
Kaplan, J., McCandlish, S., Henighan, T., et al.
\newblock \emph{Scaling Laws for Neural Language Models}.
\newblock arXiv:2001.08361, 2020.

\bibitem{hoffmann2022training}
Hoffmann, J., Borgeaud, S., Mensch, A., et al.
\newblock \emph{Training Compute-Optimal Large Language Models}.
\newblock arXiv:2203.15556, 2022.

\bibitem{vaswani2017attention}
Vaswani, A., Shazeer, N., Parmar, N., et al.
\newblock \emph{Attention Is All You Need}.
\newblock NeurIPS, 2017.

\bibitem{touvron2023llama}
Touvron, H., et al.
\newblock \emph{LLaMA: Open and Efficient Foundation Language Models}.
\newblock arXiv:2302.13971, 2023.

\bibitem{dao2022flashattention}
Dao, T., et al.
\newblock \emph{FlashAttention: Fast and Memory-Efficient Exact Attention}.
\newblock arXiv:2205.14135, 2022.

\bibitem{gu2023mamba}
Gu, A., Dao, T.
\newblock \emph{Mamba: Linear-Time Sequence Modeling with Selective State Spaces}.
\newblock arXiv:2312.00752, 2023.

\bibitem{lattner2021mlir}
Lattner, C., et al.
\newblock \emph{MLIR: Scaling Compiler Infrastructure for Domain Specific Computation}.
\newblock CGO, 2021.

\bibitem{chen2018tvm}
Chen, T., et al.
\newblock \emph{TVM: An Automated End-to-End Optimizing Compiler for Deep Learning}.
\newblock OSDI, 2018.

\end{thebibliography}

\appendix
\section{Architecture Details}

\subsection{Crates}

\begin{table}[ht]
\centering
\caption{The nine crates of the NEURAX workspace.}
\begin{tabular}{@{}l l@{}}
\toprule
\textbf{Crate} & \textbf{Responsibility} \\
\midrule
\texttt{neurax-core} & Pipeline orchestration (11 phases) \\
\texttt{neurax-ir} & 10-dialect analytical IR \\
\texttt{neurax-parser} & JSON $\rightarrow$ ModelConfig AST \\
\texttt{neurax-formulas} & Analytical cost formulas \\
\texttt{neurax-hardware-db} & GPU/CPU specifications \\
\texttt{neurax-mlir} & MLIR/LLVM 18 backend \\
\texttt{neurax-cli} & Command-line interface (\texttt{neurax} binary) \\
\texttt{neurax-service} & Actix-web HTTP API \\
\texttt{neurax-tui} & Terminal UI \\
\bottomrule
\end{tabular}
\label{tab:crates}
\end{table}

\subsection{Formula Examples}

Listing~\ref{lst:attention} shows the attention FLOPs computation as implemented in the \texttt{neurax-formulas} crate.

\begin{lstlisting}[language=RustModern,caption={Attention FLOPs (abridged, neurax-formulas/src/attention.rs).},label={lst:attention}]
// Q, K, V projections: 3 x 2 x B x S x H x H
let qkv_flops = 3.0 * (2.0 * b * s * h * h);
// QK^T scores: 2 x B x heads x S^2 x head_dim
let attn_scores_flops = 2.0 * b * heads * s * s * head_dim;
// Attention x V: 2 x B x heads x S^2 x head_dim
let attn_v_flops = 2.0 * b * heads * s * s * head_dim;
// Output projection: 2 x B x S x H x H
let out_proj_flops = 2.0 * b * s * h * h;
// Softmax: ~5 x B x heads x S^2 (exp, sum, div)
let softmax_flops = 5.0 * b * heads * s * s;
// Causal attention computes half the attention matrix
let causal_factor = if causal { 0.5 } else { 1.0 };
qkv_flops + (attn_scores_flops + attn_v_flops
    + softmax_flops) * causal_factor + out_proj_flops
\end{lstlisting}

\vspace{1em}
\noindent\textit{Preprint. For the latest version, see \url{https://github.com/rustnew/NEURAX}.}

\end{document}
